\chapter{Song Lyrics (\emph{Ci}) and After}

\section{What a \emph{ci} is}

\begin{definition}[\emph{Ci} \zht{詞}]
A \emph{cí} (``song lyric'') is a poem written to fit the melody of an existing tune. Each
tune (\zht{詞牌} \emph{cípái}, ``tune title'') fixes the number of lines, the number of
characters in each line, the tonal pattern and the rhyme positions. Writing a \emph{ci} is
called \zht{填詞} \emph{tián cí}, ``filling in words''.
\end{definition}

\begin{important}[\emph{Shi} versus \emph{ci}]
\begin{center}
\begin{tabular}{lll}
\toprule
 & \emph{Shi} \zht{詩} & \emph{Ci} \zht{詞} \\
\midrule
Line length & Uniform (5 or 7) & Irregular (1 to 11 characters), set by the tune \\
Title & Tells the occasion & Tune name (unrelated to content), sometimes plus a subtitle \\
Structure & One stanza (or none) & Usually two stanzas (\zht{上片} / \zht{下片}) \\
Origin & Literati, examinations & Entertainment quarters, songs for singing girls \\
Register & Public, moral, ``serious'' & Private, emotional, romantic (at first) \\
Height & Tang & Song (960--1279) \\
\bottomrule
\end{tabular}
\end{center}
\end{important}

\begin{remark}
The music for \emph{ci} came from the new ``banquet music'' of the Tang, influenced by
Central Asian music. The melodies are almost all lost. Writers in later centuries composed
\emph{ci} from the tonal patterns alone, as if the tune were still there, using
reference books of patterns (\zht{詞譜}).
\end{remark}

\section{Early masters: the Southern Tang}

Li Yu \zht{李煜} (937--978), last ruler of the Southern Tang, a small southern kingdom
conquered by the Song. Captured, he lived as a prisoner in the Song capital until he was
poisoned. His captivity lyrics changed \emph{ci} from party songs into personal
lament. To the tune ``The Beautiful Lady Yu'' \zht{虞美人}:

\begin{center}
\begin{tabular}{p{0.38\textwidth}p{0.52\textwidth}}
\toprule
\zht{春花秋月何時了？} & Spring flowers, autumn moon --- when will they end? \\
\zht{往事知多少。} & How much of the past do I know? \\
\zht{小樓昨夜又東風，} & Last night, in the small tower, the east wind again: \\
\zht{故國不堪回首月明中。} & I cannot bear to turn my head toward my old country in the bright moon. \\
\zht{雕闌玉砌應猶在，} & Carved railings and jade steps must be there still, \\
\zht{只是朱顏改。} & only the rosy faces have changed. \\
\zht{問君能有幾多愁？} & You ask how much sorrow can a man hold? \\
\zht{恰似一江春水向東流。} & Just like a river in spring flowing east. \\
\bottomrule
\end{tabular}
\end{center}

\textbf{Analysis.} The cycle of seasons that should be beautiful is unbearable because it
never ends. Note the line lengths (7, 5, 7, 9 / 7, 5, 7, 9): the two stanzas mirror each
other. The final simile (sorrow = an endless river in spring flood) became one of the most
quoted lines in Chinese.

\section{Two schools}

Critics later divided Song \emph{ci} into two styles:
\begin{itemize}
  \item \textbf{Delicate and restrained} (\zht{婉約} \emph{wǎnyuē}): love, longing,
        boudoir scenes, refined emotion. Liu Yong \zht{柳永}, Qin Guan, Li Qingzhao.
  \item \textbf{Heroic and unrestrained} (\zht{豪放} \emph{háofàng}): history, philosophy,
        patriotism, broad landscape. Su Shi, Xin Qiji \zht{辛棄疾}.
\end{itemize}

\begin{aside}
An anecdote: Su Shi asked a singer how his lyrics compared with Liu Yong's. The reply:
Liu Yong's should be sung by a girl of seventeen with red clappers, ``Willow bank, dawn
wind, a waning moon''; Su Shi's need a big fellow from the west with an iron clapper
singing ``The great river flows east''. The story sums up the two schools.
\end{aside}

\section{Su Shi \zht{蘇軾} (1037--1101)}

Also known as Su Dongpo \zht{蘇東坡}. Poet, essayist, calligrapher, painter, statesman,
engineer, gourmet (Dongpo pork); the most complete figure of Song culture. An opponent of
Wang Anshi's New Policies, he was imprisoned (1079, the ``Crow Terrace Poetry Case'',
accused of slandering the emperor in verse), then exiled to Huangzhou, and later to Huizhou
and finally Hainan island. He met each exile with remarkable equanimity, drawing on Daoism,
Chan Buddhism and Tao Yuanming.

\subsection{To the tune ``Prelude to Water Music'' \zht{水調歌頭} (1076)}

Preface (abridged): ``Mid-Autumn Festival of 1076: I drank happily till dawn, got very
drunk, and wrote this piece, also thinking of Ziyou.'' (Ziyou: his brother Su Zhe, whom he
had not seen for seven years.)

\begin{center}
\begin{tabular}{p{0.36\textwidth}p{0.54\textwidth}}
\toprule
\zht{明月幾時有？把酒問青天。} & When did the bright moon first appear? Cup in hand, I ask the blue sky. \\
\zht{不知天上宮闕，今夕是何年。} & I do not know, in the palaces of heaven, what year it is tonight. \\
\zht{我欲乘風歸去，} & I would ride the wind and go home there, \\
\zht{又恐瓊樓玉宇，高處不勝寒。} & but fear those towers of jasper and halls of jade --- so high, the cold too much to bear. \\
\zht{起舞弄清影，何似在人間？} & I rise and dance and play with my clear shadow: how could that compare with being among people? \\
\midrule
\zht{轉朱閣，低綺戶，照無眠。} & It turns past the red pavilion, stoops to the silk-screened window, shines on the sleepless. \\
\zht{不應有恨，何事長向別時圓？} & It should bear no grudge --- so why is it always full when people are apart? \\
\zht{人有悲歡離合，月有陰晴圓缺，} & People have sorrow and joy, parting and meeting; the moon is dim or bright, full or waning; \\
\zht{此事古難全。} & since ancient times this has never been made whole. \\
\zht{但願人長久，千里共嬋娟。} & I only wish that we may live long, and a thousand miles apart, share her beauty. \\
\bottomrule
\end{tabular}
\end{center}

\subsection{The shape of the tune}

Each bar below is one line of the lyric; its length is the number of characters, and a
filled dot marks a rhyme. The two stanzas begin differently but end with the same
cadence (6, 6, 5 / 5, 5).

\begin{center}
\begin{tikzpicture}[x=0.5cm, y=0.4cm]
  \node[font=\scriptsize, anchor=south west] at (0,0.6) {stanza 1 (\zht{上片})};
  \foreach \n/\r [count=\i] in {5/0,5/1,6/0,5/1,6/0,6/0,5/1,5/0,5/1} {
    \fill[blue!30] (0,-\i) rectangle (\n,-\i+0.75);
    \draw (0,-\i) rectangle (\n,-\i+0.75);
    \node[font=\tiny, anchor=east] at (-0.2,-\i+0.37) {\n};
    \ifnum\r=1 \fill[red!70!black] (\n+0.6,-\i+0.37) circle (0.18); \fi
  }
  \node[font=\scriptsize, anchor=south west] at (16,0.6) {stanza 2 (\zht{下片})};
  \foreach \n/\r [count=\i] in {3/0,3/0,3/1,4/0,7/1,6/0,6/0,5/1,5/0,5/1} {
    \fill[orange!35] (16,-\i) rectangle (16+\n,-\i+0.75);
    \draw (16,-\i) rectangle (16+\n,-\i+0.75);
    \node[font=\tiny, anchor=east] at (15.8,-\i+0.37) {\n};
    \ifnum\r=1 \fill[red!70!black] (16+\n+0.6,-\i+0.37) circle (0.18); \fi
  }
  \fill[red!70!black] (8,-11.6) circle (0.18);
  \node[font=\scriptsize, anchor=west] at (8.5,-11.6) {rhyme: \zht{天 年 寒 間 / 眠 圓 全 娟} (\emph{-an})};
\end{tikzpicture}
\end{center}

\subsection{Analysis}
\begin{itemize}
  \item \textbf{Stanza 1: the pull upward.} The banished-immortal theme of Li Bai (asking the
        sky, wishing to ``return'' to heaven), then a turn: heaven is too cold; better to be
        human. Dancing with his shadow recalls Li Bai's ``Drinking Alone under the
        Moon''.
  \item \textbf{Stanza 2: the moon moves.} Three three-character lines trace the moon's path
        across the house to the sleepless poet. He accuses the moon of being full at
        times of separation, then answers himself: imperfection is the nature of things,
        for people as for the moon.
  \item \textbf{Ending.} Consolation through shared perception: separated people look at the
        same moon (\zht{嬋娟}, ``the lovely one'', = the moon goddess / the moon). The
        oldest moon convention, renewed.
  \item The \emph{heroic} style here is not loudness but breadth: cosmic questions and a
        philosophical resolution, in a genre that had been about lovers.
  \item It is still sung at Mid-Autumn Festival (Teresa Teng's 1983 setting
        \zht{但願人長久} is famous).
\end{itemize}

\section{Li Qingzhao \zht{李清照} (1084--1155?)}

The greatest woman poet in the Chinese tradition. Daughter of a scholar-official, married
happily to Zhao Mingcheng, with whom she collected antiquities and inscriptions. The
Jurchen invasion of 1127 destroyed the Northern Song; the couple fled south, their
collection was lost piece by piece, and Zhao died in 1129. Her later lyrics are about
widowhood, exile and loss. She also wrote a short critical essay on the \emph{ci}
(\zht{詞論}) insisting it was a distinct art with its own musical rules, criticising
famous male writers (including Su Shi) for writing ``\emph{shi} with irregular lines''.

\subsection{To the tune ``Like a Dream'' \zht{如夢令}}

\begin{center}
\begin{tabular}{p{0.22\textwidth}p{0.3\textwidth}p{0.38\textwidth}}
\toprule
\zht{昨夜雨疏風驟，} & zuó yè yǔ shū fēng zhòu & Last night, the rain was sparse, the wind gusting; \\
\zht{濃睡不消殘酒。} & nóng shuì bù xiāo cán jiǔ & deep sleep did not clear the last of the wine. \\
\zht{試問捲簾人，} & shì wèn juǎn lián rén & I ask the maid rolling up the blind --- \\
\zht{卻道海棠依舊。} & què dào hǎitáng yī jiù & she says the crab-apple is just the same. \\
\zht{知否？知否？} & zhī fǒu? zhī fǒu? & Don't you know? Don't you know? \\
\zht{應是綠肥紅瘦。} & yīng shì lǜ féi hóng shòu & It must be the green grown fat, the red grown thin. \\
\bottomrule
\end{tabular}
\end{center}

\textbf{Analysis.} A tiny drama in 33 characters, with dialogue. The maid's careless answer
provokes the speaker's correction: after a night of wind, the blossoms (red) must have
fallen and the leaves (green) flourished. \zht{綠肥紅瘦} ``green fat, red thin'' is a
startling synaesthetic phrase (colours given body weight), and the fading flowers suggest
the speaker's own youth.

\subsection{To the tune ``Slow Sound'' \zht{聲聲慢} (opening, late period)}

\begin{center}
\begin{tabular}{p{0.36\textwidth}p{0.54\textwidth}}
\toprule
\zht{尋尋覓覓，冷冷清清，} & Searching, seeking, cold, cold, clear, clear, \\
\zht{淒淒慘慘戚戚。} & bleak, bleak, wretched, wretched, grieving, grieving. \\
\zht{乍暖還寒時候，最難將息。} & The season when it turns warm and then cold again is the hardest to rest in. \\
\zht{三杯兩盞淡酒，} & Three cups, two bowls of thin wine --- \\
\zht{怎敵他、晚來風急？} & how can they stand against the evening's sharp wind? \\
\zht{雁過也，正傷心，} & Wild geese passing: just when my heart is breaking, \\
\zht{卻是舊時相識。} & they turn out to be old acquaintances. \\
\bottomrule
\end{tabular}
\end{center}

\begin{important}[Fourteen reduplicated characters]
The opening has seven reduplicated pairs (14 characters), unprecedented in \emph{ci}. The
sound (many dental and palatal initials, \emph{x-, q-, j-, ch-}) has been compared to
sobbing. Translators struggle: repetition in English sounds childish. The geese are ``old
acquaintances'' because they once carried letters between her and her husband (the goose
messenger image from Chapter 1); now there is no one to write to.
\end{important}

\section{Later traditions and modern echoes (brief)}

\begin{itemize}
  \item \textbf{Song \emph{shi}}: more discursive, philosophical, humorous and everyday than
        Tang verse (Su Shi, Huang Tingjian, Lu You). ``Tang poetry is about feeling; Song
        poetry is about reasoning'' (a cliché, but a useful one).
  \item \textbf{Yuan \emph{qu}} \zht{曲} (``arias''): songs in a freer, colloquial language,
        used in drama and as independent lyrics (\zht{散曲}). Ma Zhiyuan's ``Autumn
        Thoughts'' (\zht{天淨沙·秋思}) is built entirely of noun phrases:
        \zht{枯藤老樹昏鴉，小橋流水人家，古道西風瘦馬。} ``Withered vines, old trees,
        crows at dusk; small bridge, flowing water, someone's home; ancient road, west
        wind, a thin horse.'' A late culmination of the imagistic juxtaposition we began
        with.
  \item \textbf{Ming and Qing}: arguments over whether to imitate the High Tang or follow
        individual ``spirit''; huge output; women's poetry flourishes (Qing anthologies of
        women poets).
  \item \textbf{Modern}: the 1917 Literary Revolution (Hu Shi \zht{胡適}) called for
        poetry in the vernacular (\zht{白話}), free verse rather than regulated forms. Yet
        classical forms never died: Mao Zedong wrote \emph{ci}, and writing regulated verse
        remains a living hobby; every Chinese schoolchild memorises Tang poems.
  \item \textbf{In English}: Pound's \emph{Cathay} (1915) and Arthur Waley's translations
        (1918 onwards) shaped modernist free verse; later, Gary Snyder (Hanshan), Kenneth
        Rexroth (Du Fu, Li Qingzhao) and David Hinton made classical Chinese poetry part of
        American poetry itself. The influence of the ``Chinese poem'' in English is a
        story of productive mistranslation.
\end{itemize}

\begin{important}[Course review: the big line]
\emph{Shijing} (four-character, \emph{xing}) $\to$ \emph{Chuci} (\zht{兮}, Qu Yuan)
$\to$ Han \emph{yuefu} and Nineteen Old Poems (five-character) $\to$ Tao Yuanming (fields
and gardens) $\to$ Tang regulated verse (Wang Wei, Li Bai, Du Fu, Bai Juyi, Li Shangyin)
$\to$ Song \emph{ci} (Li Yu, Su Shi, Li Qingzhao) $\to$ Yuan \emph{qu} $\to$ vernacular
modern poetry. Recurring threads: the moon, parting, reclusion vs service, history and
ruins, and the translator's problem of what the poem leaves unsaid.
\end{important}

\section*{Summary}
\begin{itemize}
  \item \emph{Ci}: lyrics fitted to tune patterns (\zht{詞牌}), irregular lines, two stanzas;
        ``filling in words''.
  \item Li Yu: from party song to personal lament; sorrow as a spring river.
  \item Two schools: delicate (\zht{婉約}) and heroic (\zht{豪放}).
  \item Su Shi's ``Prelude to Water Music'': moon, brother, heaven too cold, ``share her
        beauty a thousand miles apart''.
  \item Li Qingzhao: ``green fat, red thin''; fourteen reduplicated characters; widowhood
        and the geese.
  \item Afterwards: Song \emph{shi}, Yuan \emph{qu}, vernacular poetry, and translation into
        English modernism.
\end{itemize}
